\documentclass[11pt]{scrartcl}

\usepackage[top=1.5cm]{geometry}
\usepackage{url}
\usepackage{float}
\usepackage{listings}
\usepackage{xcolor}
\usepackage{graphicx}

\setlength{\parindent}{0em}
\setlength{\parskip}{0.5em}

\newcommand{\youranswerhere}{[Your answer goes here \ldots]}
\renewcommand{\thesubsection}{\arabic{subsection}}

\lstdefinestyle{dmrsql}{
  language=SQL,
  basicstyle=\small\ttfamily,
  keywordstyle=\color{magenta!75!black},
  stringstyle=\color{green!50!black},
  showspaces=false,
  showstringspaces=false,
  commentstyle=\color{gray}}

\lstdefinestyle{dmrJava}{
  language=JAVA,
  basicstyle=\small\ttfamily,
  keywordstyle=\color{magenta!75!black},
  stringstyle=\color{green!50!black},
  showspaces=false,
  showstringspaces=false,
  commentstyle=\color{gray}}

\title{
  \textbf{\large Projektaufgabe 2 } \\
  Phase 2 – Legen der Basis (2.5 P) \\
  {\large Datenmanagement jenseits von Relationen}
}

\author{
  Gruppen Nummer (A12) \\
  \large Zanotti Christian, 12528509 \\
  \large Wiesinger Martin, 11935667  \\
  \large Ottino David, 51841010
}

\begin{document}

\maketitle\thispagestyle{empty}

Dieses Reporting Template dient der Vorbereitung der Abgabe von Phase 2.

\subsection*{Datengenerator für Matrizen mit Sparsity (0.5 Punkte)}

Zeigen Sie den Code der Funktion generate() als Listing oder Screenshot. Gehen Sie
(kurz) auf die wesentlichen Aspekte ein

Die Methode "generate" wird mit den Parametern l für die Dimension und sparsity für die Dichte aufgerufen. Es wird überprüft, ob sparsity im geforderten Bereich liegt und anschließend die Dimensionen (m und n) der beiden Matrizen anhand von l ermittelt. Anschließend werden zwei 2-D-Arrays erzeugt, indem über l und m bzw. n gelooped wird und in jedem Loop ein zufälliger Integer-Wert in das Array eingefügt wird. Mithilfe von sparsity werden auch an zufälligen Positionen 0-Werte eingefügt.

\begin{lstlisting}[language=Python]
def generate(l: int, sparsity: float):

    if not (0.0 <= sparsity <= 1.0):
        raise ValueError("sparsity has to be between 0 and 1")

    m = l - 1
    n = l - 1

    A = [
        [random_value(sparsity) for _ in range(l)]
        for _ in range(m)
    ]

    B = [
        [random_value(sparsity) for _ in range(n)]
        for _ in range(l)
    ]

    return A, B

def random_value(sparsity):
    if random.random() < sparsity:
        return 0.0
    else:
        return random.uniform(-10, 10)
\end{lstlisting}

\subsection*{Import der Matrizen in das DBMS (0.5 Punkte):}

Geben Sie das Create Table Statement für Matrix A an.

\begin{lstlisting}[style=dmrsql]
CREATE TABLE IF NOT EXISTS A (
                i INT NOT NULL,
                j INT NOT NULL,
                val INT NOT NULL
            );
\end{lstlisting}

Für die Übung bereiten Sie eine Demo des Datenimports vor.


\subsection*{Wahl des Toy Beispiels (0.5 P):}

Geben Sie Matrix A und B als 2D Array an.

\begin{lstlisting}[style=dmrsql]
A:
[[1, 0, 2, 0], 
 [0, 3, 0, 4], 
 [5, 0, 0, 6]]

B:
[[7, 0, 8], 
 [0, 9, 0], 
 [1, 0, 2], 
 [0, 3, 0]]
\end{lstlisting}

Zeigen Sie die äquivalente darstellung der Matrix A und B in der Datenbank

\begin{lstlisting}[style=dmrsql]
A:
"i"	"j"	"val"
1	1	1
1	3	2
2	2	3
2	4	4
3	1	5
3	4	6

B:
"i"	"j"	"val"
1	1	7
1	3	8
2	2	9
3	1	1
3	3	2
4	2	3
\end{lstlisting}

\subsection*{Implementierung von Ansatz 0 (0.5 P):}
Zeigen Sie den Code der Matrixmultiplikation als Listing oder Screenshot. Erläutern Sie
(kurz), welche Laufzeit ihr Algorithmus hat und warum das Kriterium keinen Algorith-
mus mit sub-kubischer Laufzeit zu w¨ahlen erf¨ullt ist

\begin{lstlisting}[language=Python]
def ansatz0(A, B):
    m = len(A)
    l = len(A[0])
    n = len(B[0])

    C = [[0.0 for _ in range(n)] for _ in range(m)]

    for i in range(m):
        for j in range(n):
            for k in range(l):
                C[i][j] += A[i][k] * B[k][j]

    return C
\end{lstlisting}

Dem Algorithmus werden zwei 2-D-Arrays übergeben. Aus diesem werden die Parameter m, n und l ermittelt. Anschließend wird das Ergebnis-Array C in der richtigen Dimension initialisiert (m x n). Die Berechnung selbst entspricht dem Standard-Schema für Matrix-Multiplikation: Jede Zeile in Matrix A wird mit jeder Spalte der Matrix B multipliziert. Die Summe dieser Multiplikation ist ein Ergebnis der Matrix C. Dieser Algorithmus hat kubische Laufzeit.

\subsection*{Implementierung von Ansatz 1 (0.5 P):}
Berechnen Sie von Hand das Ergebnis $C = A \times B$ füur ihr Toy Beispiel und geben Sie es
nachfolgend an.

\begin{lstlisting}[style=dmrsql]
[[9.0, 0.0, 12.0], 
 [0.0, 39.0, 0.0], 
 [35.0, 18.0, 40.0]]
\end{lstlisting}

Zeigen Sie, dass Ihr System $C$ korrekt berechnet (z.B. als Screenshot)

\begin{lstlisting}[style=dmrsql]
(1, 1, 9)
(1, 3, 12)
(2, 2, 39)
(3, 1, 35)
(3, 2, 18)
(3, 3, 40)
\end{lstlisting}

\subsection*{Zeitmamagement}

Benötigte Zeit pro Person (nur Phase 1): \textbf{4}

\subsection*{References}

\begin{table}[H]
  \centering
  \begin{tabular}{c}
    \hline
    \textbf{Important:} Reference your information sources! \tabularnewline
    Remove this section if you use footnotes to reference your information sources. \tabularnewline
    \hline
  \end{tabular}
\end{table}

\end{document}
